Para evaluar y seleccionar el motor de \acs{ocr} más adecuado al dominio
del SCDEE, se diseñó una prueba de calidad sobre un \textit{dataset}
conformado por muestras reales de escritura manuscrita de alumnos de
la UAM y otros colaboradores. Se recopilaron un total de 156 muestras,
con textos de nombre y apellido, \acs{nia} y \acs{dni}.
Las muestras se procesaron con los dos motores
disponibles (EasyOCR y Tesseract) y se midió tanto la calidad del
reconocimiento en bruto como la precisión del sistema completo de
\textit{matching} (\cref{code:matching}) que lo consume.

\subsection{Metodología}

Se preparó un conjunto de 52 estudiantes sintéticos, cada uno con tres
imágenes de muestra (nombre y apellidos, \acs{nia} y \acs{dni}), totalizando
52 exámenes de 1 página y 156 zonas. Para cada zona, se ejecutó el reconocimiento con ambos
motores y se comparó el texto extraído con el valor esperado mediante
tres métricas: coincidencia exacta, coincidencia dentro de un margen de
tolerancia (calculada por la distancia de Levenshtein acotada por tipo de zona) y
precisión de caracteres (\textit{character accuracy}, definida como
\(1 - \text{distancia} / \max(|\text{esperado}|, |\text{obtenido}|)\)).

A continuación, se reconstruyó el flujo completo de asignación: los 52
lotes de atributos extraídos por cada motor se emparejaron contra la
lista de 52 estudiantes convocados mediante el algoritmo húngaro y la
función de \textit{scoring} por defecto (ponderaciones 20\,\% nombre,
50\,\% \acs{nia}, 30\,\% \acs{dni}). Se midió la precisión de la
asignación como el porcentaje de lotes correctamente emparejados con su
alumno.

\subsection{Resultados}

La \cref{table:ocr-results} resume los resultados de ambas fases.

\begin{table}[Resultados de las pruebas de calidad OCR]%
            {table:ocr-results}%
            {Resultados de las pruebas de reconocimiento óptico en bruto y
             de asignación completa (\textit{matching}).}
  \small
  \renewcommand{\arraystretch}{1.3}
  \begin{tabular}{l c c}
    \hline
    \textbf{Métrica} & \textbf{EasyOCR} & \textbf{Tesseract} \\
    \hline
    \multicolumn{3}{c}{\textit{Reconocimiento en bruto (156 muestras)}} \\
    \hline
    Precisión exacta          & 4{,}49\,\%  & 0{,}64\,\%  \\
    Dentro de tolerancia      & 25{,}64\,\% & 12{,}18\,\% \\
    Precisión de caracteres   & 59{,}35\,\% & 27{,}33\,\% \\
    Latencia media            & 248\,ms     & 83\,ms      \\
    \hline
    \multicolumn{3}{c}{\textit{Asignación completa (52 lotes)}} \\
    \hline
    Precisión de asignación   & \textbf{100\,\% (52/52)} & 82{,}2\,\% (37/45) \\
    Latencia media            & 55{,}33\,ms & 29{,}20\,ms \\
    \hline
  \end{tabular}
\end{table}

\subsection{Interpretación}

Los resultados muestran que, aunque EasyOCR es más lento que Tesseract
en el reconocimiento en bruto, su precisión es significativamente
superior, aunque no lo suficiente para ser funcional por sí mismo, con solo
un 4{,}5\% de exactitud. Sin embargo, al aplicar el sistema de matching
implementado en la aplicación, se consigue alcanzar el 100\,\% de precisión.

Con Tesseract, que no llega al 1\,\% de precisión, también se consigue un
resultado aceptable con un 82{,}2\,\%, cometiendo 8 errores
de asignación.

La combinación de EasyOCR con la función de \textit{scoring} tolerante a
confusiones de \acs{ocr} y el algoritmo húngaro resulta, por tanto, la
configuración más fiable para el escenario de identificación de alumnos
del SCDEE, y es la que se utiliza como configuración por defecto del
sistema. Además tan solo añade al rededor de 55ms de latencia por página
asociada.

Los resultados detallados de cada muestra y cada motor se recogen en
\cref{AP:TEST-RESULTS}.
